\section*{Bab \ref{ch:g1:landscapes-around-us} --- Bentang Alam di Sekitar Kita}

\begin{solution}{exo:g1:landscapes-around-us:1}
\omterm{def:g1:landscapes-around-us:landform}{Gunung}, \omterm{def:g1:landscapes-around-us:landform}{bukit}, \omterm{def:g1:landscapes-around-us:landform}{lembah}, \omterm{def:g1:landscapes-around-us:landform}{dataran}.
\end{solution}

\begin{solution}{exo:g1:landscapes-around-us:2}
\omterm{def:g1:landscapes-around-us:landform}{Gunung}: \omterm{def:g1:landscapes-around-us:landform}{gunung} lebih tinggi dan lebih curam daripada \omterm{def:g1:landscapes-around-us:landform}{bukit}.
\end{solution}

\begin{solution}{exo:g1:landscapes-around-us:3}
Alami: \omterm{def:g1:landscapes-around-us:water}{sungai}, hutan, pantai. Buatan: jalan, jembatan, rumah.
\end{solution}

\begin{solution}{exo:g1:landscapes-around-us:4}
Tempat daratan bertemu laut.
\end{solution}

\begin{solution}{exo:g1:landscapes-around-us:5}
Misalnya: yang alami, \omterm{def:g1:landscapes-around-us:landform}{gunung-gunung}, hutan di lerengnya, atau \omterm{def:g1:landscapes-around-us:water}{sungai}
kecilnya; yang buatan, sebuah rumah atau gereja desa itu.
\end{solution}

\begin{solution}{exo:g1:landscapes-around-us:6}
\omterm{def:g1:landscapes-around-us:landform}{Gunung}, yang puncaknya bersalju. \omterm{def:g1:landscapes-around-us:water}{Sungai} mengalir di dasar \omterm{def:g1:landscapes-around-us:landform}{lembah}, di antara
\omterm{def:g1:landscapes-around-us:landform}{gunung-gunung} dan \omterm{def:g1:landscapes-around-us:landform}{bukit}.
\end{solution}

\begin{solution}{exo:g1:landscapes-around-us:7}
\omterm{def:g1:landscapes-around-us:water}{Sungai} mengalir dari tanah yang lebih tinggi (\omterm{def:g1:landscapes-around-us:landform}{bukit} dan \omterm{def:g1:landscapes-around-us:landform}{gunung}) turun ke
\omterm{def:g1:landscapes-around-us:water}{danau} atau ke laut. Air selalu mengalir ke bawah, jadi air tidak bisa
mengalir naik ke puncak \omterm{def:g1:landscapes-around-us:landform}{bukit}.
\end{solution}

\begin{solution}{exo:g1:landscapes-around-us:8}
Pohon itu dekat dengan kita, di latar depan; \omterm{def:g1:landscapes-around-us:landform}{bukit} itu jauh, di latar
belakang. Benda-benda yang jauh tampak lebih kecil daripada benda-benda
yang dekat.
\end{solution}

\begin{solution}{exo:g1:landscapes-around-us:9}
Latar depan: rumput tinggi dan beberapa semak di bagian bawah gambar.
Latar tengah: \omterm{def:g1:landscapes-around-us:water}{sungai} yang berkelok-kelok melintasi \omterm{def:g1:landscapes-around-us:landform}{dataran} dengan
pohon-pohon di tepinya, ladang, sebuah pertanian, dan sebuah jalan.
Latar belakang: tanah datar yang terus membentang sampai ke cakrawala, dan
langit.
\end{solution}

\begin{solution}{exo:g1:landscapes-around-us:10}
Karena manusia yang membuatnya: mereka membuka lahan, membajaknya, lalu
menabur benih gandum. Tanpa manusia, di sana akan ada rumput, semak, atau
hutan, bukan ladang gandum.
\end{solution}

\begin{solution}{exo:g1:landscapes-around-us:11}
Tidak berubah: \omterm{def:g1:landscapes-around-us:landform}{bukit} dan \omterm{def:g1:landscapes-around-us:water}{sungai}. Ditambahkan oleh manusia: ladang, lebih
banyak rumah, jalan, dan jembatan (dan sebagian besar hutannya sudah
ditebang).
\end{solution}
